\section{Discussion}
\label{sec:discussion}

\paragraph{What the results establish}
The offline results establish two things of different kinds. The first concerns the data. On the shipped files, the order in which the release lists the evidence places answer-bearing passages among the first three less often than chance would on three of four datasets, and a reordering that needs no labels repairs this. No agent takes part in this measurement, so it holds for any backbone. The second concerns the code. Under replies of known content, the release departs from the specification of the paper in a share of the cases, and MADKE-X does not. How large that share is with a real model depends on how the model formats its replies, and this study does not know it. A model that always writes \texttt{[Yes]} and never writes a list in one bracket would trigger neither finding C3 nor C4. The swapped summariser prompt, the consistency formula and the handling of errors are different: they do not depend on the format of any reply.

\paragraph{Why the order matters for a debate}
In MADKE an agent chooses at most three passages from a listing that opens with five to ten Wikipedia passages, while the short Google snippets that most often contain the answer string come last. Language models use long contexts unevenly \cite{liu2024lost}, so an agent is more likely to pick early passages that help less, and the debate then rests on weaker evidence. The RRF order puts an answer-bearing passage among the first three for \VGainMin{} to \VGainMax{} percentage points more questions, from the question and the retrieval rank alone. Whether agents turn the better listing into better answers is what the live comparison measures; the offline result shows that the opportunity exists. The comparison of orderings also carries a caution. Simply interleaving the two sources achieves most of the gain, and on NQ it achieves more than RRF at three passages. The merit of the fusion is that it holds up on every dataset. It is not the best choice on each.

\paragraph{The price of the specified pool}
The pool that the MADKE paper specifies makes the listing \VTokRedMean\% shorter than the default of the release, and it removes the only answer-bearing passage for \VCoverLossMin{} to \VCoverLossMax\% of the questions. A faithful implementation has to follow the specification, and the pool size remains configurable through $w$ and $g$ in Eq.~\eqref{eq:pool}. On HotpotQA and 2WikiMultihopQA about half of the questions or fewer have an answer-bearing passage in either pool. Retrieval, and not debate, limits these datasets, and iterative retrieval \cite{trivedi2023ircot} is the natural next step for them.

\paragraph{Traceability as a property of the system}
For a question-answering system in use, showing why an answer was given matters apart from accuracy. The trace of MADKE-X records what the pipeline did and is in that sense faithful to the process \cite{jacovi2020faithful}: every answer can be inspected, failures are counted, and a comparison can be recomputed from the stored records. The lexical support check is the weakest part, and deliberately so. It detects that the answer string is missing from the cited passages. It does not establish that a passage entails the answer, and it rejects correct paraphrases. A learned attribution model \cite{rashkin2023attribution,gao2023alce} could replace it without touching the rest of the pipeline. The verification step built on the check inherits its limits. It can never replace a supported answer and costs at most one call. Whether supported answers are more often correct than unsupported ones can be read directly from the live records, which store the status next to the gold answers.

\paragraph{What is not shown}
This study does not show that MADKE-X answers more questions correctly than MADKE. That claim rests on the live runs of Table~\ref{tab:live}. Nor does it show that the published results of MADKE are wrong. The findings concern the public code, and the published shares fit an evaluation of all 500 questions.
